\documentclass[11pt, a4paper]{article}

% bfw-style.tex — gemeinsame Style-Definitionen für alle Handouts/Aufgabenhefte
% Voraussetzung: Kompilation mit xelatex (wegen fontspec)

\usepackage[ngerman]{babel}
% Babel-ngerman macht " zu einem aktiven Shorthand (z.B. für "a -> ä). In unseren
% Texten verwenden wir " als normales Zeichen — daher Shorthand komplett ausschalten.
\shorthandoff{"}
\usepackage{geometry}
\geometry{a4paper, top=2.4cm, bottom=2.4cm, left=2.3cm, right=2.3cm,
          headheight=15pt, headsep=0.7cm, footskip=1.1cm}

\usepackage{fontspec}
% Fallback-Kette: Source Sans Pro -> Calibri -> Segoe UI -> Latin Modern Sans
% (Latin Modern ist immer mit MiKTeX vorhanden)
\IfFontExistsTF{Source Sans Pro}{%
  \setmainfont{Source Sans Pro}[Ligatures=TeX]%
}{\IfFontExistsTF{Calibri}{%
  \setmainfont{Calibri}[Ligatures=TeX]%
}{\IfFontExistsTF{Segoe UI}{%
  \setmainfont{Segoe UI}[Ligatures=TeX]%
}{%
  \setmainfont{Latin Modern Sans}[Ligatures=TeX]%
}}}
\IfFontExistsTF{Source Code Pro}{%
  \setmonofont{Source Code Pro}[Scale=0.92]%
}{\IfFontExistsTF{Consolas}{%
  \setmonofont{Consolas}[Scale=0.92]%
}{%
  \setmonofont{Latin Modern Mono}[Scale=0.92]%
}}

% Gesperrte Versalien für Labels/Titelbalken (Kapitälchen-Ersatz, funktioniert
% mit jeder OpenType-Schrift — Latin Modern Sans hat keine echten Kapitälchen)
\newcommand{\bfwlabelfont}{\footnotesize\bfseries\addfontfeatures{LetterSpace=8.0}}

\usepackage{xcolor}
% Farbpalette: hell, freundlich, modern
\definecolor{bfwPrimary}{HTML}{0EA5A4}    % Petrol/Türkis - Primär
\definecolor{bfwPrimaryDark}{HTML}{0F766E} % dunkler Petrol für Headings
\definecolor{bfwAccent}{HTML}{F59E0B}     % Amber - Akzent
\definecolor{bfwSuccess}{HTML}{10B981}    % Grün - Erfolg / Lösung
\definecolor{bfwWarn}{HTML}{F97316}       % Orange - Achtung
\definecolor{bfwInk}{HTML}{0F172A}        % fast Schwarz
\definecolor{bfwInkSoft}{HTML}{475569}    % gedämpftes Grau
\definecolor{bfwBgSoft}{HTML}{F8FAFC}     % off-white
\definecolor{bfwBgTeal}{HTML}{ECFEFF}     % helles Türkis
\definecolor{bfwBgAmber}{HTML}{FEF3C7}    % helles Gelb
\definecolor{bfwBgRose}{HTML}{FFE4E6}     % helles Rosé für Eselsbrücken
\definecolor{bfwTabHead}{HTML}{E4F3F2}    % zarte Kopfzeilen-Hinterlegung für Tabellen

\usepackage[colorlinks=true, linkcolor=bfwPrimaryDark, urlcolor=bfwPrimaryDark]{hyperref}
\usepackage{lastpage}
\usepackage{enumitem}
\setlist[itemize]{leftmargin=*, itemsep=2pt, topsep=2pt}
\setlist[enumerate]{leftmargin=*, itemsep=2pt, topsep=2pt}

\usepackage[most]{tcolorbox}
\usepackage{booktabs}
\usepackage{colortbl}
\usepackage{tikz}
\usetikzlibrary{decorations.pathmorphing, positioning, shapes.geometric, arrows.meta}

% ---------------------------------------------------------------------------
% Einheitliches Tabellen-Design
%   - etwas mehr Zeilenluft, dezente graue Linien (wirkt auch mit \hline ruhiger)
%   - Kopfzeile einer Tabelle bei Bedarf zart hinterlegen: \rowcolor{bfwTabHead}
% ---------------------------------------------------------------------------
\renewcommand{\arraystretch}{1.25}
\arrayrulecolor{bfwInkSoft!50}
\setlength{\heavyrulewidth}{1pt}
\setlength{\lightrulewidth}{0.5pt}

% ---------------------------------------------------------------------------
% Boxen — klare farbige Titelbalken mit gesperrten Versal-Labels
% (Titel bleibt per Option überschreibbar: \begin{merksatz}[title={...}])
% ---------------------------------------------------------------------------
\tcbset{bfwbox/.style={%
  enhanced, breakable,
  boxrule=0.5pt, arc=2.5pt,
  left=10pt, right=10pt, top=8pt, bottom=8pt,
  toptitle=4pt, bottomtitle=4pt,
  titlerule=0pt,
  fonttitle=\bfwlabelfont,
  coltitle=white
}}

% Merkkästchen — wichtige Definition / Kernpunkt
\newtcolorbox{merksatz}[1][]{%
  bfwbox,
  colback=bfwBgTeal,
  colframe=bfwPrimaryDark,
  colbacktitle=bfwPrimaryDark,
  title={MERKE},
  #1
}

% Eselsbrücke — Mnemoniks
\newtcolorbox{eselsbruecke}[1][]{%
  bfwbox,
  colback=bfwBgRose,
  colframe=bfwAccent!85!black,
  colbacktitle=bfwAccent!85!black,
  title={ESELSBRÜCKE},
  #1
}

% Beispiel-Box
\newtcolorbox{beispiel}[1][]{%
  bfwbox,
  colback=bfwBgSoft,
  colframe=bfwInkSoft,
  colbacktitle=bfwInkSoft,
  title={BEISPIEL},
  #1
}

% Lösungs-Box (für Aufgabenhefte)
\newtcolorbox{loesung}[1][]{%
  bfwbox,
  colback=bfwSuccess!7,
  colframe=bfwSuccess!65!black,
  colbacktitle=bfwSuccess!65!black,
  title={LÖSUNG},
  #1
}

% Achtung-Box — typische Fehler / Stolperfallen
\newtcolorbox{achtung}[1][]{%
  bfwbox,
  colback=bfwWarn!6,
  colframe=bfwWarn!85!black,
  colbacktitle=bfwWarn!85!black,
  title={ACHTUNG},
  #1
}

% Formel-Box — für zentrale Formeln (groß, ruhig, ohne Titelbalken)
\newtcolorbox{formelbox}[1][]{%
  enhanced,
  colback=bfwBgTeal,
  colframe=bfwPrimaryDark,
  boxrule=1pt, arc=3pt,
  left=10pt, right=10pt, top=6pt, bottom=6pt,
  #1
}

% Glossar-Eintrag
\newcommand{\glossareintrag}[2]{%
  \par\noindent\textbf{\color{bfwPrimaryDark}#1} \enspace #2\par\smallskip
}

% ---------------------------------------------------------------------------
% Abschnitts-Design: farbiges Nummern-Quadrat + feine Linie unter \section,
% dezent farbige \subsection/\subsubsection
% ---------------------------------------------------------------------------
\usepackage{titlesec}
\newcommand{\bfwSecBadge}[1]{%
  \begin{tikzpicture}[baseline={([yshift=-0.62em]n.center)}]
    \node[fill=bfwPrimaryDark, text=white, font=\normalsize\bfseries,
          minimum width=1.55em, minimum height=1.55em, inner sep=1pt,
          rounded corners=1.5pt] (n) {#1};
  \end{tikzpicture}}
\titleformat{\section}
  {\Large\bfseries\color{bfwPrimaryDark}}
  {\bfwSecBadge{\thesection}}{0.6em}{}
  [\vspace{-0.2em}{\color{bfwPrimary!60}\titlerule[0.8pt]}]
\titleformat{\subsection}
  {\large\bfseries\color{bfwPrimaryDark}}
  {\textcolor{bfwPrimary}{\thesubsection}}{0.5em}{}
\titleformat{\subsubsection}
  {\normalsize\bfseries\color{bfwInk}}
  {\textcolor{bfwPrimary}{\thesubsubsection}}{0.4em}{}
\titlespacing*{\section}{0pt}{1.6em}{1em}
\titlespacing*{\subsection}{0pt}{1.2em}{0.5em}

% TikZ-Stil "skizzenhaft" — etwas weicher als Rough.js, aber stabil
\tikzset{
  roughbox/.style = {
    draw=bfwPrimaryDark, thick, fill=bfwBgTeal,
    rounded corners=4pt, inner sep=6pt
  },
  roughaccent/.style = {
    draw=bfwAccent, thick, fill=bfwBgAmber,
    rounded corners=4pt, inner sep=6pt
  },
  roughline/.style = {
    draw=bfwInkSoft, thick, -{Stealth[length=2.5mm]}
  }
}

% Bullet-Symbole (bewusst kein FontAwesome — vermeidet Font-Installations-Probleme)
\newcommand{\faLightbulb}{$\bullet$}
\newcommand{\faBookmark}{$\bullet$}

% ---------------------------------------------------------------------------
% Kopf-/Fußzeile
%   Kopf links:  Wortlogo „Wirtschaftsfachwirt"
%   Kopf rechts: Dokument-Kurztext, gesetzt via \kopfzeile{...}
%   Fuß links:   © Can Siebert · 2026 — Fuß rechts: Seite X von Y
%   Titelseite (Seite 1) bleibt ohne Kopf-/Fußzeile (\handouttitel setzt
%   \thispagestyle{empty}).
% ---------------------------------------------------------------------------
\usepackage{fancyhdr}
\newcommand{\bfwKopfzeileText}{}
\newcommand{\kopfzeile}[1]{\renewcommand{\bfwKopfzeileText}{#1}}
\pagestyle{fancy}
\fancyhf{}
\fancyhead[L]{\small\bfseries\color{bfwPrimaryDark}Wirtschaftsfachwirt}
\fancyhead[R]{\small\color{bfwInkSoft}\bfwKopfzeileText}
\renewcommand{\headrulewidth}{0.6pt}
\renewcommand{\headrule}{{\color{bfwPrimary}\hrule width\headwidth height 0.6pt}}
\fancyfoot[L]{\small\color{bfwInkSoft}\copyright\ Can Siebert\,\textperiodcentered\,2026}
\fancyfoot[R]{\small\color{bfwInkSoft}Seite \thepage\ von \pageref*{LastPage}}
% Auch \thispagestyle{plain} (z.B. durch \maketitle) soll leer bleiben:
\fancypagestyle{plain}{\fancyhf{}\renewcommand{\headrulewidth}{0pt}\renewcommand{\headrule}{}}

% ---------------------------------------------------------------------------
% \handouttitel{Obertitel}{Untertitel}{Kurs-Label}{Termin-Zeile}
%   Einheitlicher professioneller Titelblock: Dachzeile, farbige Headline,
%   Untertitel, farbige Trennlinie und dezente Meta-Zeile (Kurs/Termin/Dozent).
%   Ersetzt \title/\maketitle. Direkt nach \begin{document} aufrufen.
% ---------------------------------------------------------------------------
\newcommand{\handouttitel}[4]{%
  \thispagestyle{empty}%
  \begingroup
  \setlength{\parindent}{0pt}%
  {\bfwlabelfont\color{bfwPrimary}WIRTSCHAFTSFACHWIRT (IHK)\par}
  \vspace{0.55em}
  {\huge\bfseries\color{bfwPrimaryDark}#1\par}
  \vspace{0.5em}
  {\large\color{bfwInkSoft}#2\par}
  \vspace{0.9em}
  {\color{bfwPrimary}\rule{\linewidth}{2pt}}\par
  \vspace{0.55em}
  {\small\color{bfwInkSoft}#3\ \,\textperiodcentered\,\ #4\ \,\textperiodcentered\,\ Dozent: Can Siebert\par}
  \endgroup
  \vspace{1.4em}%
}


\kopfzeile{Aufgabenheft BWL Lektion 2}

\begin{document}
\handouttitel{Aufgabenheft Lektion~2 (BWL)}%
  {Betriebliche Funktionen II: Rechnungswesen, Finanzierung, Controlling \& Personal --- Übungen im IHK-Klausurformat}%
  {1.2 Betriebswirtschaftliche Grundlagen}%
  {Selbstlernmaterial zur Lektion vom Sa, 12.09.2026}


\begin{merksatz}[title={\bfseries So nutzen Sie dieses Aufgabenheft}]
Bearbeiten Sie alle Aufgaben zunächst \emph{ohne} in die Lösungen zu schauen. Schreiben
Sie Ihre Antworten in vollständigen Sätzen --- die IHK-Prüfung verlangt formulierte
Begründungen, nicht bloße Stichworte. Beachten Sie die \emph{Operatoren} (nennen,
erläutern, ordnen zu, beurteilen, analysieren) --- sie geben den geforderten
Antwort-Tiefegrad vor. Erst danach öffnen Sie den Lösungsteil ab
Seite~\pageref{loesungen} und vergleichen. Ergänzend: Übungsband Betriebswirtschaft,
Kapitel~1 (Übungen 25--44).
\end{merksatz}

\tableofcontents
\vspace{1em}
\hrule

\section{Aufgaben}

\subsection*{Aufgabe~1 --- Aufgaben und Teilgebiete des Rechnungswesens (12 Punkte)}

\begin{enumerate}[label=(\alph*)]
  \item \textbf{Nennen und erläutern} Sie die vier Aufgaben des betrieblichen
    Rechnungswesens. \hfill (8~P.)
  \item \textbf{Nennen} Sie die vier Teilgebiete des Rechnungswesens. \hfill (4~P.)
\end{enumerate}

\subsection*{Aufgabe~2 --- Internes vs. externes Rechnungswesen (8 Punkte)}

\begin{enumerate}[label=(\alph*)]
  \item \textbf{Unterscheiden} Sie internes und externes Rechnungswesen nach Teilgebieten,
    Adressaten und Rechtsgrundlagen. \hfill (6~P.)
  \item \textbf{Ordnen} Sie die folgenden Instrumente jeweils dem internen oder externen
    Rechnungswesen \textbf{zu}: die Bilanz zum 31.12.; die Kalkulation des Verkaufspreises
    für ein neues Produkt. \hfill (2~P.)
\end{enumerate}

\subsection*{Aufgabe~3 --- Investitionsarten (10 Punkte)}

Investitionen sind Vorgänge der Kapitalverwendung.

\begin{enumerate}[label=(\alph*)]
  \item \textbf{Ordnen} Sie die folgenden Beispiele jeweils der Sachinvestition,
    der Finanzinvestition oder der immateriellen Investition \textbf{zu}:
    (1) Kauf einer Immobilie, (2) Kauf einer Lizenz für eine neue Kühltechnologie,
    (3) Erwerb von Anteilen an einem anderen Haushaltsgerätehersteller,
    (4) Kurse zur Gesundheitsvorsorge für die Mitarbeiter,
    (5) Anschaffung einer Filteranlage für die Abgase. \hfill (5~P.)
  \item \textbf{Unterscheiden} Sie Erst-, Ersatz- und Erweiterungsinvestition und
    \textbf{ordnen} Sie folgenden Fall \textbf{zu}: Eine defekte Produktionsanlage wird
    durch eine modernere, leistungsfähigere Anlage ersetzt. \hfill (5~P.)
\end{enumerate}

\subsection*{Aufgabe~4 --- Zuordnungsaufgabe Finanzierungsarten (12 Punkte)}

\begin{enumerate}[label=(\alph*)]
  \item \textbf{Ordnen} Sie die folgenden Finanzierungsvorgänge in die Matrix aus
    Außen-/Innenfinanzierung und Eigen-/Fremdkapital ein (bei Kapitalfreisetzung genügt
    die Angabe „Innenfinanzierung aus Kapitalfreisetzung``):
    \begin{enumerate}[label=\arabic*.]
      \item Es tritt ein neuer Gesellschafter in das Unternehmen ein.
      \item Gewinne werden einbehalten und für Investitionen genutzt.
      \item Das Unternehmen nimmt einen Kredit auf.
      \item Die Abschreibungsbeträge für eine Maschine werden zum Kauf einer neuen
        Maschine eingesetzt.
      \item Das Unternehmen least eine Maschine.
      \item Das Unternehmen bildet Pensionsrückstellungen und nutzt die Mittel bis zur
        Fälligkeit.
      \item Das Unternehmen verkauft Grundstücke, um Maschinen zu beschaffen.
    \end{enumerate} \hfill (7~P.)
  \item \textbf{Erläutern} Sie die Begriffe Außenfinanzierung, Innenfinanzierung,
    Eigenfinanzierung und Fremdfinanzierung. \hfill (5~P.)
\end{enumerate}

\subsection*{Aufgabe~5 --- Eigenkapital vs. Fremdkapital (8 Punkte)}

\textbf{Ordnen} Sie die folgenden Merkmale dem Eigenkapital oder dem Fremdkapital
\textbf{zu}: \hfill (je 1~P.)

\begin{enumerate}[label=(\alph*)]
  \item erfolgsabhängige Gewinnbeteiligung
  \item Haftung erfolgt in Höhe der Einlage
  \item vorrangiger Anspruch auf Rückzahlung der Forderung
  \item steht unbefristet zur Verfügung
  \item Kapitalgeber ist Gläubiger
  \item Kapitalgeber ist grundsätzlich zur Mitbestimmung berechtigt
  \item Zinsen sind steuerlich absetzbar
  \item Es besteht ein Beteiligungsverhältnis
\end{enumerate}

\subsection*{Aufgabe~6 --- Leasing, Factoring \& goldene Finanzierungsregel (10 Punkte)}

\begin{enumerate}[label=(\alph*)]
  \item \textbf{Erläutern} Sie je einen Vorteil und einen Nachteil des Leasings aus Sicht
    des Leasingnehmers und \textbf{beurteilen} Sie die Aussage: „Leasing ist eine Form der
    Selbstfinanzierung.`` \hfill (4~P.)
  \item \textbf{Erläutern} Sie zwei Wirkungen des Factorings auf die Liquidität des
    Unternehmens. \hfill (4~P.)
  \item \textbf{Formulieren} Sie die goldene Finanzierungsregel. \hfill (2~P.)
\end{enumerate}

\subsection*{Aufgabe~7 --- Controlling (12 Punkte)}

Die Süd IT-Systeme e.\,K. hat eine neue Stelle im Bereich Controlling geschaffen. Ein
Mitarbeiter der internen Revision meint, Controlling sei überflüssig, da es schon die
Kontrolle durch die Revision gebe.

\begin{enumerate}[label=(\alph*)]
  \item \textbf{Grenzen} Sie Controlling und Kontrolle voneinander \textbf{ab} und nehmen
    Sie damit Stellung zur Aussage des Mitarbeiters. \hfill (4~P.)
  \item \textbf{Beschreiben} Sie die vier Aufgabenbereiche des Controllings im
    Controlling-Kreislauf (Planung, Kontrolle, Steuerung, Information). \hfill (4~P.)
  \item \textbf{Unterscheiden} Sie strategisches und operatives Controlling nach
    Leitfrage, Zeithorizont und Zielgrößen. \hfill (4~P.)
\end{enumerate}

\subsection*{Aufgabe~8 --- Personalwirtschaft (10 Punkte)}

\begin{enumerate}[label=(\alph*)]
  \item \textbf{Nennen und erläutern} Sie die fünf Kernaufgaben der Personalwirtschaft
    (Personalplanung, -beschaffung, -entwicklung, -einsatz, -freisetzung). \hfill (5~P.)
  \item \textbf{Ordnen} Sie die folgenden Tätigkeiten jeweils einem Aufgabenbereich der
    Personalwirtschaft \textbf{zu}: Führung der Personalakten; Durchführung eines
    Assessmentcenters; Konzeption von Weiterbildungsangeboten; Analyse der Personalkosten;
    Erstellung von Arbeitsplatzbeschreibungen. \hfill (3~P.)
  \item \textbf{Unterscheiden} Sie wirtschaftliche und soziale Ziele der
    Personalwirtschaft mit je einem Beispiel. \hfill (2~P.)
\end{enumerate}

\subsection*{Aufgabe~9 --- Fallaufgabe: Zielkonflikte zwischen Funktionsbereichen (14 Punkte)}

Bei der Süd Möbelwerke GmbH treffen in der Jahresplanung folgende Forderungen aufeinander:
Der \emph{Vertrieb} verlangt mehr Produktvarianten und kürzere Lieferzeiten, um Kunden zu
gewinnen. Die \emph{Produktion} fordert Standardisierung und große Fertigungslose, um
Rüstzeiten und Stückkosten zu senken. Der \emph{Einkauf} will große Bestellmengen ordern,
um Mengenrabatte zu nutzen. Die \emph{Finanzabteilung} mahnt eine geringe Kapitalbindung
im Lager und schnelle Zahlungseingänge an.

\begin{enumerate}[label=(\alph*)]
  \item \textbf{Erklären} Sie, was unter einem Zielkonflikt zwischen betrieblichen
    Funktionsbereichen zu verstehen ist. \hfill (2~P.)
  \item \textbf{Analysieren} Sie zwei der geschilderten Zielkonflikte: beteiligte
    Bereiche, jeweilige Interessen und Kern des Konflikts. \hfill (6~P.)
  \item \textbf{Schlagen} Sie der Unternehmensleitung zwei Koordinationsinstrumente
    \textbf{vor}, mit denen die Konflikte entschärft werden können, und \textbf{begründen}
    Sie Ihre Vorschläge. \hfill (4~P.)
  \item \textbf{Erläutern} Sie, welchen Beitrag das Denkmodell der Wertschöpfungskette
    zur Lösung solcher Konflikte leistet. \hfill (2~P.)
\end{enumerate}

\newpage
\section{Lösungen}\label{loesungen}

\begin{loesung}[title={Lösung Aufgabe~1 --- Aufgaben und Teilgebiete des Rechnungswesens}]
\textbf{(a)} \emph{Dokumentation:} lückenlose, planmäßige Aufzeichnung aller
Geschäftsvorfälle anhand von Belegen (keine Buchung ohne Beleg).
\emph{Information (Rechenschaftslegung):} Unterrichtung interner und externer Adressaten
(Geschäftsleitung, Gesellschafter, Gläubiger, Finanzamt) über die Lage des Unternehmens.
\emph{Kontrolle:} Überwachung von Wirtschaftlichkeit, Rentabilität und Liquidität, z.\,B.
durch Soll-Ist-Vergleiche und Nachkalkulation.
\emph{Disposition:} Bereitstellung des Zahlenmaterials als Grundlage für Planungen und
unternehmerische Entscheidungen.

\textbf{(b)} (1) Buchführung (Finanzbuchhaltung mit Jahresabschluss),
(2) Kosten- und Leistungsrechnung, (3) Statistik (Vergleichsrechnung),
(4) Planungsrechnung (Vorschaurechnung).
\end{loesung}

\begin{loesung}[title={Lösung Aufgabe~2 --- Internes vs. externes Rechnungswesen}]
\textbf{(a)} \emph{Externes ReWe:} Buchführung und Jahresabschluss (Bilanz, GuV);
Adressaten sind außenstehende Gruppen (Gläubiger, Gesellschafter, Finanzamt,
Öffentlichkeit); gesetzlich verpflichtend nach HGB, AO und Steuergesetzen;
vergangenheitsorientiert. \emph{Internes ReWe:} Kosten- und Leistungsrechnung, Statistik,
Planungsrechnung; Adressat ist die Unternehmensführung; keine gesetzlichen Vorschriften,
freiwillig und zweckorientiert; gegenwarts- und zukunftsorientiert.

\textbf{(b)} Bilanz $\rightarrow$ externes Rechnungswesen; Verkaufspreiskalkulation
$\rightarrow$ internes Rechnungswesen (KLR).
\end{loesung}

\begin{loesung}[title={Lösung Aufgabe~3 --- Investitionsarten}]
\textbf{(a)} (1) Immobilie: \emph{Sachinvestition}. (2) Lizenz: \emph{immaterielle
Investition}. (3) Anteile an anderem Hersteller: \emph{Finanzinvestition}.
(4) Gesundheitskurse: \emph{immaterielle Investition} (Investition in das Humankapital).
(5) Filteranlage: \emph{Sachinvestition}.

\textbf{(b)} \emph{Erstinvestition:} erstmalige Ausstattung bei Gründung (auch
Gründungsinvestition). \emph{Ersatzinvestition:} Ersatz einer verbrauchten oder defekten
Anlage durch eine gleichartige. \emph{Erweiterungsinvestition:} Ausbau der Kapazität über
den bisherigen Bestand hinaus. Der geschilderte Fall ist eine \emph{Ersatzinvestition} ---
da die neue Anlage leistungsfähiger ist, verbindet sie sich mit einem
Rationalisierungseffekt (Rationalisierungsinvestition).
\end{loesung}

\begin{loesung}[title={Lösung Aufgabe~4 --- Finanzierungsarten}]
\textbf{(a)}
(1) neuer Gesellschafter: \emph{Außenfinanzierung, Eigenkapital}
(Einlagen-/Beteiligungsfinanzierung).
(2) Gewinneinbehaltung: \emph{Innenfinanzierung, Eigenkapital} (Selbstfinanzierung).
(3) Kreditaufnahme: \emph{Außenfinanzierung, Fremdkapital} (Kreditfinanzierung).
(4) Abschreibungsgegenwerte: \emph{Innenfinanzierung aus Kapitalfreisetzung}
(Abschreibungsfinanzierung).
(5) Leasing: \emph{Außenfinanzierung, Fremdkapital} (Sonderform der Kreditfinanzierung /
Kreditsubstitut).
(6) Pensionsrückstellungen: \emph{Innenfinanzierung, Fremdkapital}
(Rückstellungsfinanzierung).
(7) Grundstücksverkauf: \emph{Innenfinanzierung aus Kapitalfreisetzung}
(Vermögensumschichtung).

\textbf{(b)} \emph{Außenfinanzierung:} Dem Unternehmen werden zusätzliche Mittel von
außen zugeführt (durch Eigentümer oder Gläubiger). \emph{Innenfinanzierung:} Die Mittel
stammen aus dem betrieblichen Umsatzprozess und werden im Betrieb belassen.
\emph{Eigenfinanzierung:} Zuführung bzw. Bildung von Eigenkapital (Einlagen, Beteiligungen,
einbehaltene Gewinne). \emph{Fremdfinanzierung:} Zuführung von Fremdkapital --- es
entsteht ein Schuldverhältnis gegenüber Gläubigern (Kredite, Anleihen, Rückstellungen).
\end{loesung}

\begin{loesung}[title={Lösung Aufgabe~5 --- Eigenkapital vs. Fremdkapital}]
(a) Eigenkapital --- der Gewinnanteil hängt vom Erfolg ab.
(b) Eigenkapital --- der Gesellschafter haftet mindestens mit seiner Einlage.
(c) Fremdkapital --- Gläubiger werden in der Insolvenz vorrangig bedient.
(d) Eigenkapital --- es steht grundsätzlich unbefristet zur Verfügung.
(e) Fremdkapital --- es besteht ein Schuldverhältnis.
(f) Eigenkapital --- Eigentümer haben Mitsprache- und Kontrollrechte.
(g) Fremdkapital --- Fremdkapitalzinsen mindern als Aufwand den steuerpflichtigen Gewinn.
(h) Eigenkapital --- der Kapitalgeber ist am Unternehmen beteiligt.
\end{loesung}

\begin{loesung}[title={Lösung Aufgabe~6 --- Leasing, Factoring \& goldene Finanzierungsregel}]
\textbf{(a)} \emph{Vorteil:} Bei der Anschaffung ist nur ein geringer Finanzmittelbedarf
nötig --- die Liquidität wird geschont, der Investitionsspielraum wird größer.
\emph{Nachteil:} Die Summe der Leasingraten ist in der Regel höher als der Kaufpreis;
beim Operate-Leasing erwirbt der Leasingnehmer zudem kein Eigentum am Gegenstand.
\emph{Beurteilung:} Die Aussage ist falsch --- Leasing ist keine Selbstfinanzierung,
sondern eine Sonderform der Außen-/Fremdfinanzierung (Kreditsubstitut); Selbstfinanzierung
ist die Einbehaltung von Gewinnen.

\textbf{(b)} Durch den laufenden Verkauf der Forderungen fließt der Gegenwert sofort zu:
(1) Liquiditätsengpässe durch lange Zahlungsfristen der Kunden werden vermieden;
(2) mit der gewonnenen Liquidität können Lieferantenskonti genutzt und enge Kreditrahmen
der Banken geschont werden. (Zusätzlich möglich: Übernahme des Ausfallrisikos
--- Delkrederefunktion --- und Auslagerung des Debitorenmanagements.)

\textbf{(c)} Goldene Finanzierungsregel (Fristenkongruenz): Die Dauer der
Kapitalüberlassung soll der Dauer der Kapitalbindung entsprechen --- langfristig
gebundenes Vermögen (Anlagevermögen) ist langfristig zu finanzieren, kurzfristig
gebundenes Vermögen darf kurzfristig finanziert werden.
\end{loesung}

\begin{loesung}[title={Lösung Aufgabe~7 --- Controlling}]
\textbf{(a)} \emph{Kontrolle} ist der vergangenheitsbezogene Soll-Ist-Vergleich ---
sie stellt Abweichungen fest. \emph{Controlling} ist eine umfassende
Management-Service-Funktion: Es umfasst neben der Kontrolle auch Planung,
Informationsversorgung und die zukunftsgerichtete Steuerung der Unternehmensprozesse.
Der Einwand des Mitarbeiters greift daher zu kurz: Die interne Revision prüft
vergangenheitsorientiert die Ordnungsmäßigkeit; das Controlling unterstützt die Führung
aktiv und zukunftsorientiert bei Planung und Steuerung --- beide ersetzen einander nicht.

\textbf{(b)} \emph{Planung:} Ziele, Maßnahmen und Budgets festlegen (Sollgrößen).
\emph{Kontrolle:} Soll-Ist-Vergleich und Abweichungsanalyse.
\emph{Steuerung:} Gegensteuerungsmaßnahmen einleiten und Pläne anpassen.
\emph{Information:} entscheidungsrelevante Informationen beschaffen, aufbereiten und der
Führung bereitstellen (Berichtswesen). Die vier Bereiche bilden einen fortlaufenden
Kreislauf.

\textbf{(c)} \emph{Strategisches Controlling:} Leitfrage „Tun wir die richtigen
Dinge?``; langfristiger Zeithorizont; Zielgrößen: Erfolgspotenziale und
Existenzsicherung (Chancen/Risiken, Stärken/Schwächen).
\emph{Operatives Controlling:} Leitfrage „Tun wir die Dinge richtig?``; kurz- bis
mittelfristig; Zielgrößen: Gewinn, Rentabilität, Liquidität und Wirtschaftlichkeit
(Budgets, Soll-Ist-Vergleiche, Kennzahlen).
\end{loesung}

\begin{loesung}[title={Lösung Aufgabe~8 --- Personalwirtschaft}]
\textbf{(a)} \emph{Personalplanung:} künftigen Personalbedarf nach Menge, Qualifikation,
Zeit und Ort ermitteln. \emph{Personalbeschaffung:} offene Stellen intern (Versetzung)
oder extern (Anzeigen, Jobportale) besetzen. \emph{Personalentwicklung:} Mitarbeiter durch
Aus- und Weiterbildung fördern und Potenziale erschließen. \emph{Personaleinsatz:}
Mitarbeiter den Arbeitsplätzen anforderungs- und eignungsgerecht zuordnen.
\emph{Personalfreisetzung:} personelle Überkapazitäten möglichst sozialverträglich
abbauen (Fluktuation, Altersteilzeit, zuletzt Kündigung).

\textbf{(b)} Personalakten $\rightarrow$ Personalverwaltung; Assessmentcenter
$\rightarrow$ Personalbeschaffung (Auswahl); Weiterbildungskonzeption $\rightarrow$
Personalentwicklung; Personalkostenanalyse $\rightarrow$ Personalcontrolling;
Arbeitsplatzbeschreibungen $\rightarrow$ Personaleinsatz.

\textbf{(c)} \emph{Wirtschaftliche Ziele:} optimale Versorgung mit qualifizierten
Mitarbeitern bei optimalen Personalkosten (z.\,B. hohe Arbeitsproduktivität).
\emph{Soziale Ziele:} Zufriedenheit und Motivation der Mitarbeiter (z.\,B. gerechte
Entlohnung, Vereinbarkeit von Familie und Beruf).
\end{loesung}

\begin{loesung}[title={Lösung Aufgabe~9 --- Zielkonflikte zwischen Funktionsbereichen}]
\textbf{(a)} Ein Zielkonflikt liegt vor, wenn die Ziele einzelner Funktionsbereiche
einander widersprechen: Die bessere Erreichung des einen Bereichsziels verschlechtert
die Erreichung des anderen --- obwohl beide Bereiche zum gemeinsamen Unternehmensziel
beitragen sollen.

\textbf{(b)} Beispiel 1 --- \emph{Vertrieb vs. Produktion:} Der Vertrieb will
Produktvielfalt und kurze Lieferzeiten (Kundenorientierung, Umsatz), die Produktion will
Standardisierung und große Lose (geringe Rüstzeiten, niedrige Stückkosten). Kern:
Variantenvielfalt erhöht Komplexität und Kosten der Fertigung.
Beispiel 2 --- \emph{Einkauf vs. Finanzen/Lager:} Der Einkauf will große Bestellmengen
(Mengenrabatte, Versorgungssicherheit), die Finanzabteilung will geringe Kapitalbindung
und niedrige Lagerkosten. Kern: Der Rabattvorteil wird durch Lager- und Zinskosten
(gebundenes Kapital) erkauft. (Ebenfalls möglich: Vertrieb --- lange Zahlungsziele vs.
Finanzen --- schnelle Zahlungseingänge.)

\textbf{(c)} Mögliche Koordinationsinstrumente: (1) eine \emph{abgestimmte
Unternehmensplanung mit Budgets und Zielvereinbarungen} durch das Controlling --- die
Bereichsziele werden aus den Unternehmenszielen abgeleitet und aufeinander abgestimmt;
(2) \emph{bereichsübergreifende Gremien/Prozesse} (z.\,B. Sales-and-Operations-Planning,
gemeinsame Sortiments- und Bestandsplanung), in denen die Bereiche Zielgrößen wie
Variantenzahl, Losgrößen und Lagerreichweiten gemeinsam festlegen. (Auch denkbar:
Kennzahlensysteme, Verrechnungspreise, klare Priorisierung durch die Geschäftsleitung.)

\textbf{(d)} Die Wertschöpfungskette macht sichtbar, dass alle Funktionen --- primäre wie
unterstützende --- Glieder \emph{eines} Leistungsprozesses sind und nur gemeinsam
Wertschöpfung erzeugen. Sie lenkt den Blick von Bereichsinteressen auf die Schnittstellen
und den Gesamtbeitrag zum Kundennutzen und Unternehmensziel.
\end{loesung}

\vspace{1em}
\hrule
\vspace{0.8em}
\noindent\textsf{\small Weiterüben: Übungsband Betriebswirtschaft Kap.~1, Übungen 25--44
(mit Lösungsteil) und das interaktive Quiz dieser Lektion. Nächste Lektion (Sa, 26.09.2026):
Existenzgründung \& Rechtsformen.}

\end{document}
